% Generated draft: requires mathematical and source review.
\documentclass[11pt]{article}
\usepackage[margin=25mm]{geometry}
\usepackage{fontspec}
\setmainfont{Latin Modern Roman}
\usepackage{amsmath,amssymb,mathtools,unicode-math}
\setmathfont{Latin Modern Math}
\usepackage{enumitem,needspace,xurl,hyperref}
\setlength{\parindent}{0pt}
\setlength{\parskip}{5pt}
\setlength{\emergencystretch}{3em}
\newcommand{\N}{\mathbb{N}}
\newcommand{\Z}{\mathbb{Z}}
\newcommand{\Q}{\mathbb{Q}}
\newcommand{\R}{\mathbb{R}}
\newcommand{\C}{\mathbb{C}}
\newcommand{\F}{\mathbb{F}}
\newcommand{\A}{\mathbb{A}}
\newcommand{\PP}{\mathbb{P}}
\newcommand{\E}{\mathbb{E}}
\newcommand{\eps}{\varepsilon}
\newcommand{\dd}{\,\mathrm{d}}
\begin{document}
\section*{UnsolvedMath ID: 30001635}
\textbf{Source problem code:} catalogue entry
\subsection{Definitions and mathematical statement}
The supplied record asks about the Johnson cokernels
\[
C_g(k):=\mathfrak h_g(k)\big/\tau_k\bigl(\mathcal M_{g,1}[k+1]\bigr).
\]
For a homomorphism \(f:G\to A\) with abelian target \(A\), its cokernel is the quotient group \(A/\operatorname{im}(f)\). Thus, if \(\mathfrak h_g(k)\) is an abelian group and \(\tau_k:\mathcal M_{g,1}[k+1]\to\mathfrak h_g(k)\) is the intended homomorphism, the displayed expression has this meaning. These are conditions for interpreting the notation, not hypotheses explicitly provided by the record.

\textbf{Problem.} Determine the structure and topological meaning of \(C_g(k)\), with the notation and parameter ranges intended in the cited Oberwolfach problem.

The supplied text does not define \(\mathfrak h_g(k)\), \(\mathcal M_{g,1}[k+1]\), or \(\tau_k\), specify the allowed values of \(g\) and \(k\), or identify the coefficient ring. In particular, the filtration indexing \([k+1]\) is retained exactly and is not replaced by a presumed convention. The intended algebraic structure to be determined and the meaning of a topological interpretation are also unspecified. Consequently, this is a faithful statement of the supplied research question, but a fully specified mathematical formulation requires the missing source conventions.
\subsection{Short English statement}
Determine the algebraic structure and topological interpretation of the displayed Johnson cokernels. The source record leaves the parameter ranges, coefficients, filtration conventions, and maps undefined, so these require clarification before mathematical review.
\subsection{Sources}
\begin{itemize}
\item Local UnsolvedMath registry, \texttt{data/unsolvedmath/problems.json}, permanent ID \(30001635\), source problem code \texttt{catalogue entry}, title ``Structure and Topology of Johnson Cokernels.'' The record attributes the question to the Oberwolfach Reports collection; no individual proposer is supplied. The proposed year is recorded as 2010.
\item Supplied source citation: \emph{Teichmüller Theory} (2011), Oberwolfach Reports, pp.\ 3085--3157, DOI \texttt{10.4171/owr/2010/53}. \url{https://doi.org/10.4171/owr/2010/53}. This is also the supplied literature source labeled ``original Oberwolfach problem report.''
\item Supplied literature source labeled ``torsion in Johnson cokernels'': \url{https://academic.oup.com/imrn/article/2025/10/rnaf121/8139876}.
\item The supplied literature assessment, checked on \texttt{2026-08-21}, labels the problem ``partially solved'' and reports computations in several degrees and known torsion obstructions, while describing the full cokernel structure and topological meaning as open. This is supplied dated curation, not an independently verified literature assessment.
\item The registry labels the statement \texttt{corrected\_verified} and reports that its mathematical content was checked against the official Oberwolfach report, with notation, context, and question form normalized. This verification claim is supplied metadata and has not been independently verified here.
\end{itemize}
\end{document}
